\documentclass[11pt,a4paper]{article}
\usepackage{amsmath,amssymb,amsfonts}
\usepackage{graphicx}
\usepackage[margin=1in]{geometry}
\usepackage{longtable}

\title{Supplementary Material:\\Blind Exploration of Fundamental Constants}

\author{S-wind Algebraic System\\Valhalla Project}

\date{\today}

\begin{document}

\maketitle

\tableofcontents
\newpage

%==============================================================================
\section{Hypothesis Space Definition}
%==============================================================================

To ensure auditability, we explicitly define the hypothesis space $\mathcal{H}$ within which S-wind conducted its blind search.

\subsection{Allowed Operations}

The search was restricted to:

\begin{enumerate}
    \item \textbf{Input domain}: Integer values $D \in \{1, 2, \ldots, 10\}$
    \item \textbf{Polynomial class}: Univariate polynomials of degree $\leq 4$ with integer coefficients $|c_i| \leq 10$
    \item \textbf{Basic operations}: $\{+, -, \times, \div, \sqrt{\cdot}, \ln(\cdot)\}$
    \item \textbf{Comparison metric}: Relative error $\epsilon = |P(D) - C| / C$ for constant $C$
    \item \textbf{Match threshold}: $\epsilon < 0.001$ (0.1\%)
\end{enumerate}

\subsection{Explicitly Forbidden}

The following were \textbf{not allowed}:

\begin{itemize}
    \item Arbitrary scaling constants (no free parameters to fit)
    \item Rounding of physical constants to integers
    \item Pre-selected anchor points
    \item Assumed recurrence forms
    \item Post-hoc adjustment of match threshold
\end{itemize}

\subsection{Hypothesis Space Size}

For degree-4 polynomials with coefficients in $[-10, 10]$:
\begin{equation}
|\mathcal{H}_{\text{poly}}| = 21^5 = 4,084,101 \text{ candidate polynomials}
\end{equation}

Input domain size: $|D| = 10$ values.

Total evaluations: $\sim 4 \times 10^7$.

%==============================================================================
\section{Complete Physical Constants Table}
%==============================================================================

All 21 constants used in the blind exploration, with full precision:

\begin{longtable}{|l|l|r|l|}
\hline
\textbf{Symbol} & \textbf{Name} & \textbf{Value} & \textbf{Source} \\
\hline
\endfirsthead
\hline
\textbf{Symbol} & \textbf{Name} & \textbf{Value} & \textbf{Source} \\
\hline
\endhead
$\alpha^{-1}$ & Fine structure inverse & 137.035999084 & CODATA 2022 \\
$m_p/m_e$ & Proton/electron mass & 1836.15267343 & CODATA 2022 \\
$m_n/m_e$ & Neutron/electron mass & 1838.68366173 & CODATA 2022 \\
$m_\mu/m_e$ & Muon/electron mass & 206.7682830 & CODATA 2022 \\
$m_\tau/m_e$ & Tau/electron mass & 3477.23 & PDG 2024 \\
$m_W/m_e$ & W boson/electron mass & 157298 & PDG 2024 \\
$m_Z/m_e$ & Z boson/electron mass & 178450 & PDG 2024 \\
$m_H/m_e$ & Higgs/electron mass & 245034 & PDG 2024 \\
$m_t/m_e$ & Top quark/electron mass & 338600 & PDG 2024 \\
$m_b/m_e$ & Bottom quark/electron mass & 8194 & PDG 2024 \\
$m_c/m_e$ & Charm quark/electron mass & 2488 & PDG 2024 \\
$m_s/m_e$ & Strange quark/electron mass & 183 & PDG 2024 \\
$m_u/m_e$ & Up quark/electron mass & 4.3 & PDG 2024 \\
$m_d/m_e$ & Down quark/electron mass & 9.4 & PDG 2024 \\
$\alpha_s(M_Z)$ & Strong coupling at $M_Z$ & 0.1180 & PDG 2024 \\
$\sin^2\theta_W$ & Weinberg angle & 0.23122 & CODATA 2022 \\
$G_F$ & Fermi constant (GeV$^{-2}$) & $1.1663788 \times 10^{-5}$ & CODATA 2022 \\
$\Omega_\Lambda$ & Dark energy density & 0.6889 & Planck 2018 \\
$\Omega_m$ & Matter density & 0.3111 & Planck 2018 \\
$H_0$ & Hubble constant (km/s/Mpc) & 67.4 & Planck 2018 \\
$\ln(M_{Pl}/m_e)$ & Planck/electron log ratio & 51.528 & Derived \\
\hline
\end{longtable}

%==============================================================================
\section{Search Audit Log}
%==============================================================================

\subsection{Polynomial Search Results}

The following polynomials were evaluated at $D \in \{2,3,4,5,6,7,8\}$ and compared against all 21 constants:

\begin{verbatim}
SEARCH LOG - S-wind v1.1.1
==========================
Search space: degree-4 polynomials, |coeff| <= 10
Input range: D in {2..8}
Match threshold: epsilon < 0.001

MATCHES FOUND:
--------------
Polynomial: D^4 - 2*D^3 + 2*D + 1
  D=4: P(4) = 137
  Match: alpha^{-1} = 137.036
  Relative error: 0.00026 (0.026%)
  Status: MATCH (epsilon < 0.001)

NO OTHER MATCHES at epsilon < 0.001 threshold.

NEAR MISSES (0.001 < epsilon < 0.01):
-------------------------------------
  D^4 - D^3 - D^2 + 2 at D=5: 397 vs m_tau/m_e (error 0.88%)
  D^3 + D^2 - D + 3 at D=6: 213 vs m_mu/m_e (error 3.0%)
  [truncated - no significant patterns]

SEARCH STATISTICS:
------------------
Polynomials evaluated: 4,084,101
Total comparisons: 4,084,101 * 7 * 21 = 600,362,847
Matches at epsilon < 0.001: 1
Matches at epsilon < 0.01: 0 additional
Runtime: 3.2 seconds (parallelized)
\end{verbatim}

\subsection{Ratio Search Results}

\begin{verbatim}
RATIO SEARCH LOG
================
Searched: all pairs (C_i, C_j) for i != j
Hypothesis: C_i / C_j = p/q for small integers p,q in [1,20]

APPROXIMATE MATCHES (error < 1%):
---------------------------------
alpha^{-1} / (m_mu/m_e) = 0.663 ~ 2/3 (error 0.59%)
(m_p/m_e) / (m_n/m_e) = 0.9986 ~ 1 (error 0.14%)
(m_W/m_e) / (m_Z/m_e) = 0.8815 ~ 7/8 (error 0.74%)

EXACT MATCHES: None found.
\end{verbatim}

%==============================================================================
\section{Classification Criteria}
%==============================================================================

Results are classified into three categories with explicit criteria:

\subsection{Category A: Proved Identity}

\textbf{Criterion}: Algebraically verifiable statement that holds for all values in domain by symbolic computation.

\textbf{Example}: $P(4) = 4^4 - 2 \cdot 4^3 + 2 \cdot 4 + 1 = 137$

This is arithmetic, not a conjecture. Verification: exact rational arithmetic.

\subsection{Category B: Observed Pattern (Tested Range)}

\textbf{Criterion}: Empirical observation that holds for all tested inputs, but has not been proved for the general case.

\textbf{Example}: $P(D)$ is prime for $D \in \{2, 4, 6, 8\}$.

\textbf{Tested range}: Even integers $D \in [2, 8]$.

\textbf{Not tested}: $D \in \{10, 12, 14, \ldots\}$.

\textbf{Status}: Observation, not theorem.

\subsection{Category C: Falsified Under Model}

\textbf{Criterion}: Conjecture that has been shown to fail within a specified theoretical framework.

\textbf{Example}: ``$\exists \mu: \alpha^{-1}(\mu) = 137$'' is falsified under Standard Model QED.

\textbf{Evidence}: QED running is monotonically decreasing; below $m_e$, electron contribution decouples.

\textbf{Caveat}: May hold in extensions of SM.

%==============================================================================
\section{Permutation Test Framework}
%==============================================================================

To assess whether the P(4)=137 match is statistically significant, we outline a permutation test protocol (to be executed as future work).

\subsection{Protocol}

\begin{enumerate}
    \item \textbf{Null hypothesis}: The match between $P(D)$ and $\alpha^{-1}$ is a chance occurrence given the search space.
    
    \item \textbf{Permutation procedure}:
    \begin{itemize}
        \item Generate $N = 10000$ random permutations of the 21 constant labels
        \item For each permutation, run the same blind search
        \item Count how often a match with $\epsilon < 0.00026$ (or better) is found
    \end{itemize}
    
    \item \textbf{p-value}: Fraction of permutations achieving equal or better match.
    
    \item \textbf{Multiple comparison correction}: Bonferroni or FDR as appropriate.
\end{enumerate}

\subsection{Expected Outcome}

Given $|\mathcal{H}_{\text{poly}}| \approx 4 \times 10^6$ polynomials, 7 input values, and 21 constants:

\begin{itemize}
    \item Total comparisons: $\sim 6 \times 10^8$
    \item Expected matches at $\epsilon < 0.00026$ by chance: $\sim 6 \times 10^8 \times 0.00052 \approx 300,000$
\end{itemize}

\textbf{Critical observation}: The large search space means many spurious matches are expected. The significance of $P(4) = 137$ must be evaluated \textit{relative to search effort}, not in isolation.

\subsection{What Would Constitute Evidence}

\begin{itemize}
    \item If $P(4) = 137$ match survives with $p < 0.01$ after multiple comparison correction: \textbf{Suggestive}
    \item If no other polynomial at $D \in \{2,\ldots,8\}$ matches any constant at $\epsilon < 0.001$: \textbf{Unique within tested space}
    \item If the polynomial has independent mathematical interest (e.g., appears in number theory): \textbf{Additional weight}
\end{itemize}

%==============================================================================
\section{RG Running Details}
%==============================================================================

\subsection{QED Running of $\alpha^{-1}(\mu)$}

The one-loop QED beta function:
\begin{equation}
\mu \frac{d\alpha}{d\mu} = \frac{2\alpha^2}{3\pi} \sum_f Q_f^2 \theta(\mu - m_f)
\end{equation}

For $\mu > m_e$, the electron contributes $Q_e^2 = 1$.

\subsection{Behavior Below $m_e$}

Below the electron mass threshold:
\begin{itemize}
    \item The electron contribution \textbf{decouples} from the RG running
    \item Residual running from non-perturbative QCD effects is negligible for $\alpha$
    \item Scheme-dependent low-energy effective couplings exist but do not significantly alter $\alpha^{-1}$
\end{itemize}

\textbf{Conclusion}: $\alpha^{-1}(\mu)$ cannot reach 137 through SM running; the value at $\mu = m_e$ is already $137.036$.

\subsection{Full Standard Model Running}

Including all SM particles:
\begin{equation}
\alpha^{-1}(M_Z) \approx 127.9
\end{equation}

Running \textit{decreases} $\alpha^{-1}$ at higher scales, moving \textit{away} from 137.

%==============================================================================
\section{Auto-Generated Models: Caveats}
%==============================================================================

The following models were auto-generated by the S-wind Conjecture Container System. They are presented as \textbf{hypotheses for downstream testing}, not as physically grounded theories.

\subsection{DimensionalCoupling Model}

\textbf{Hypothesis}: $\alpha^{-1}(D) = D^4 - 2D^3 + 2D + 1$

\textbf{Status}: Generated conjecture, not derived from physics.

\textbf{Known issue}: Effective spacetime dimension in compactified theories does not simply substitute into coupling constants. A proper treatment requires:
\begin{itemize}
    \item Kaluza-Klein mode summation
    \item Volume factors from compact dimensions
    \item Running in higher-dimensional gauge theory
\end{itemize}

\textbf{Value}: The model provides a falsifiable prediction that can be tested against explicit string compactifications.

\subsection{Prime Dimension Conjecture}

\textbf{Hypothesis}: Stable physics requires $P(D)$ to be prime.

\textbf{Status}: Observation-based speculation.

\textbf{Known issue}: No known physics mechanism links primality to stability.

\textbf{Value}: Identifies an unexpected number-theoretic pattern for further investigation.

%==============================================================================
\section{Reproducibility}
%==============================================================================

\subsection{Code Availability}

All search code is available in the S-wind repository:
\begin{verbatim}
/home/jouly/AI/Valhalla/s-wind-sim/src/bin/unified_blind_exploration.rs
/home/jouly/AI/Valhalla/s-wind-sim/src/bin/conjecture_explorer.rs
/home/jouly/AI/Valhalla/s-wind-sim/src/conjecture.rs
\end{verbatim}

\subsection{Execution}

\begin{verbatim}
cd s-wind-sim
cargo run --bin unified_blind_exploration
cargo run --bin conjecture_explorer
\end{verbatim}

\subsection{Output Verification}

All numerical results can be independently verified using arbitrary-precision arithmetic (S-wind EPR module) or standard computer algebra systems.

\end{document}
